\chapter{Software Development}
\label{chap:software-development}

%% PROPOSAL: Sections 5.1-5.3 hidden for proposal, showing only Schedule (5.4)

%% \section{Software Development Methodology}
%% \label{section:software-development-methodology}
%% 
%% The project uses an Agile-inspired iterative approach adapted for a two-person team over one semester. Development is organized into two-week sprints aligned with the phases in the project schedule. Each sprint ends with a working, demonstrable increment. A shared Kanban board (Notion) tracks tasks, and GitHub is used for version control with a feature-branch workflow. Weekly check-ins with the project advisor provide feedback and course correction.
%% 
%% Given the data dependency at the start of the project (มคอ.2 documents from faculty), the first two weeks are allocated to data collection and pipeline setup, with feature development beginning only once the core data layer is operational.

%% \section{Technology Stack}
%% \label{section:technology-stack}

%% \begin{itemize}[leftmargin=80pt]
%%     \item \textbf{Frontend:} Next.js 14 (App Router), Tailwind CSS, Shadcn/UI. Chosen for familiarity, strong TypeScript support, and built-in i18n routing for Thai/English.
%% 
%%     \item \textbf{Backend:} Python FastAPI. Chosen for Python ecosystem access (PyMuPDF for PDF processing, sentence-transformers, neo4j driver) and async performance.
%% 
%%     \item \textbf{Vector database:} Supabase (PostgreSQL + pgvector extension). Stores มคอ.2 document embeddings, precomputed O*NET occupation interest profiles, user profiles, TCAS structured data, and the interaction log (program saves, card views, TCAS guide views, chatbot queries with weights and timestamps). A background job periodically recomputes collaborative scores from the interaction log. Supabase Auth handles authentication.
%% 
%%     \item \textbf{Graph database:} Neo4j on Railway. Stores the PLO--Skill--Career knowledge graph. Queried via Cypher from the FastAPI backend.
%% 
%%     \item \textbf{AI:} Gemini 2.5 Flash Lite via OpenRouter for generation, local multilingual-E5 for embeddings, Gemini text mode for structured extraction.
%% 
%%     \item \textbf{Data processing:} PyMuPDF for page text extraction, Typhoon OCR via opentyphoon.ai for low-yield pages, isolated full scanned-PDF OCR through the configured OCR model when needed, local multilingual-E5 embeddings, Gemini text-mode structured extraction, and custom Python chunking.
%% 
%%     \item \textbf{Thai NLP:} PyThaiNLP --- Thai text normalization used in TCAS program name entity resolution (fuzzy matching with sequence similarity scoring).
%% 
%%     \item \textbf{Deployment:} Vercel (frontend), Railway (FastAPI backend + Neo4j).
%% \end{itemize}
%% 
%% \section{Coding Standards}
%% \label{section:coding-standards}
%% 
%% \begin{itemize}[leftmargin=80pt]
%%     \item \textbf{Python:} PEP 8, type hints on all function signatures, docstrings for all public functions.
%% 
%%     \item \textbf{TypeScript:} ESLint with Next.js recommended config, Prettier for formatting.
%% 
%%     \item \textbf{Git:} Conventional commits format (feat/fix/chore/docs). Feature branches merged via pull request with at least one review.
%% 
%%     \item \textbf{Environment:} All secrets managed via .env files, never committed. Supabase and Gemini API keys stored as environment variables.
%% \end{itemize}

\section{Schedule}
\label{section:schedule}

Development is organized into four phases across April~2026 and January--March~2027. P1~Foundation runs April~1--12, P2~Core~AI POC runs April~6--30, P3~MVP Feature Completion runs January~2027, and P4~Polish \& Evaluation runs February--March~2027. By the end of April~2026, the current proof of concept demonstrates the data ingestion pipeline, grounded RAG chatbot, and Program Explorer over indexed curriculum content. Recommendation, explanation, and PLO/profile visualisation remain MVP features to be completed and evaluated during P3--P4.

\begin{table}[htbp]
\centering
\footnotesize
\renewcommand{\arraystretch}{1.4}
\caption{Project Schedule}
\label{tab:schedule}
\begin{tabularx}{\textwidth}{|L{3cm}|L{3.2cm}|Y|}
\hline
\textbf{Period} & \textbf{Phase} & \textbf{Deliverables} \\
\hline
Apr 1--5, 2026 & P1 · Foundation — Set up Repo & Repository setup, project structure, development environment \\
\hline
Apr 1--5, 2026 & P1 · Foundation — Data Collection & Receive มคอ.2 and TCAS admission data from KU faculty, download O*NET dataset \\
\hline
Apr 6--12, 2026 & P1 · Foundation — Data Ingestion & PDF extraction, chunking, embedding into Supabase pgvector; graph schema exploration for Phase~2 \\
\hline
Apr 6--30, 2026 & P2 · Core AI POC — RAG Chatbot & RAG pipeline over indexed curriculum and TCAS/fee data, grounded answer generation, source citations \\
\hline
Apr 6--30, 2026 & P2 · Core AI POC — Program Explorer & Semantic search and browsing over indexed program content and metadata using Supabase/pgvector \\
\hline
Apr 6--30, 2026 & P2 · Core AI POC — Demo Integration & Working proof-of-concept demo combining ingestion output, Program Explorer, and RAG chatbot \\
\hline
Jan 1--31, 2027 & P3 · MVP Features — Interest Discovery & RIASEC elicitation flow, profile scoring, confidence adjustment, and profile summary \\
\hline
Jan 1--31, 2027 & P3 · MVP Features — Program Detail Pages & Program detail pages, static PLO/profile chart, and available TCAS information \\
\hline
Jan 1--31, 2027 & P3 · Phase~2 Prototype — Portfolio Coach & Optional prototype UI only; excluded from MVP pass/fail evaluation \\
\hline
Jan 31, 2027 & P3 · MVP Features — Complete Recommendation & Complete simplified Pipeline~B2 recommendation, evidence-grounded explanation, and Program Explorer refinement \\
\hline
Feb 1--28, 2027 & P4 · Polish \& Eval — UI/UX Polish & UI/UX polish, Thai/English language switching, mobile-first refinement \\
\hline
Feb 1--28, 2027 & P4 · Polish \& Eval — Cross-feature Integration & Chatbot + Explorer + Recommendation + Program Detail integration \\
\hline
Feb 1--28, 2027 & P4 · Polish \& Eval — User Testing & RAGAS plus custom TCAS/fee regression, MRR/NDCG test set, SUS user testing with 10--15 students \\
\hline
Mar 1--31, 2027 & P4 · Polish \& Eval — Performance Testing & Performance benchmarking against all success metric targets \\
\hline
Mar 1--31, 2027 & P4 · Polish \& Eval — Bug Fixing \& Optimization & Bug fixes, optimization, stability improvements \\
\hline
Mar 1--31, 2027 & P4 · Polish \& Eval — Final Documentation & Final report, demo preparation, project documentation \\
\hline
\end{tabularx}
\end{table}

%% PROPOSAL: Progress Tracking Report hidden
%% \section{Progress Tracking Report}
%% \label{section:progress-tracking-report}
%% 
%% \begin{figure}[h]
%%     \centering
%%     \fbox{\parbox{0.9\textwidth}{\centering\vspace{1em}
%%     \textit{[GitHub contribution graph and sprint burndown charts to be included here as the project progresses.]}
%%     \vspace{1em}}}
%%     \caption{Progress Tracking}
%%     \label{fig:progress-tracking}
%% \end{figure}
